\documentclass[11pt]{article}
\usepackage[T1]{fontenc}
\usepackage{lmodern}
\usepackage{amsmath,amssymb,amsthm,mathtools}
\usepackage[margin=1in]{geometry}
\usepackage{microtype}
\usepackage[colorlinks=true,linkcolor=blue,citecolor=blue,urlcolor=blue]{hyperref}
\hypersetup{pdftitle={Quadratic Lyapunov obstructions and instantaneous vorticity growth for the three-dimensional Navier--Stokes equations},pdfauthor={}}
\newtheorem{theorem}{Theorem}[section]
\newtheorem{proposition}[theorem]{Proposition}
\newtheorem{lemma}[theorem]{Lemma}
\newtheorem{corollary}[theorem]{Corollary}
\theoremstyle{remark}
\newtheorem{remark}[theorem]{Remark}
\newcommand{\R}{\mathbb R}
\newcommand{\C}{\mathbb C}
\newcommand{\Pj}{\mathbb P}
\newcommand{\N}{\mathcal N}
\newcommand{\curl}{\operatorname{curl}}
\newcommand{\diver}{\operatorname{div}}
\newcommand{\supp}{\operatorname{supp}}
\newcommand{\diag}{\operatorname{diag}}
\newcommand{\mean}{\operatorname{mean}}
\newcommand{\norm}[1]{\left\lVert#1\right\rVert}
\newcommand{\abs}[1]{\left\lvert#1\right\rvert}
\newcommand{\ip}[2]{\left\langle#1,#2\right\rangle}
\newcommand{\dd}{\,\mathrm d}
\allowdisplaybreaks
\setlength{\emergencystretch}{2em}
\title{Quadratic Lyapunov obstructions and instantaneous vorticity growth\newline for the three-dimensional Navier--Stokes equations}
\author{}
\date{6 September 2026}
\begin{document}
\maketitle
\begin{abstract}
We establish two obstructions to particular a priori arguments for the
original unforced incompressible Navier--Stokes equation on $\R^3$.
First, no fixed symmetric quadratic form continuous on a finite-order
Sobolev space can both dominate the squared $L^3$ norm and be nonincreasing
on every local classical solution. No translation, rotation, locality,
Fourier-multiplier, or scaling assumption is imposed on the form. Amplitude
testing, averaging over translations and orthogonal transformations, and a
localized three-component Fourier interaction reduce its averaged form to
kinetic energy. Second, for each positive viscosity we construct smooth
compactly supported initial data with identical energy, enstrophy, and
maximum vorticity, whose corresponding solutions are all globally smooth,
but whose initial maximum-vorticity growth rates are unbounded. The initial
supports lie in one compact set, and every positive-order spatial vorticity
derivative vanishes at the exhibited maximizing point. This rules out a
finite instantaneous maximum-growth bound based solely on those scalar
values and that local vorticity jet. Both results concern precisely specified
classes of estimates; neither proves nor disproves arbitrary-data global
regularity. The qualitative quadratic-energy obstruction is established
background; the scope of the whole-space formulation and of the matched-scalar
construction is stated explicitly.
\end{abstract}

\section{Introduction and statements}
\label{sec:introduction}
We consider
\begin{equation}
 \partial_tu+(u\cdot\nabla)u+\nabla p=\nu\Delta u,
 \qquad \diver u=0,\qquad \nu>0,
 \label{eq:NS}
\end{equation}
on $\R^3$, with solenoidal Schwartz initial data. All solutions used below
are solutions of \eqref{eq:NS}; no averaged nonlinearity or forcing is
introduced. The first result is an impossibility theorem for a class of
quadratic functionals. The second is a construction within a class of
\emph{globally smooth} solutions. Neither statement requires a singular
Navier--Stokes solution.

Let $H^m_\sigma$ be the real closed subspace of divergence-free vector fields
in $H^m(\R^3;\R^3)$. Norms without a specified domain are on $\R^3$.

\begin{theorem}[Quadratic obstruction]\label{thm:quadratic}
Fix an integer $m\ge1$ and a viscosity $\nu>0$. There is no continuous
symmetric real bilinear form $B$ on $H^m_\sigma$ for which both of the
following hold:
\begin{enumerate}
\item There is $c>0$ such that $B(v,v)\ge c\norm v_3^2$ for every
$v\in H^m_\sigma$.
\item For the local classical solution of \eqref{eq:NS} from every
solenoidal Schwartz datum, $t\mapsto B(u(t),u(t))$ is nonincreasing on some
initial existence interval.
\end{enumerate}
The form and the coercivity constant may depend on $\nu$, but not on the
datum or time. No spatial symmetry is a hypothesis on $B$.
\end{theorem}

Continuity means $|B(v,w)|\le C\norm v_{H^m}\norm w_{H^m}$ for a finite
constant $C$. This permits arbitrary matrix Fourier weights and spatially
inhomogeneous kernels satisfying that bound. It excludes neither nonlinear
functionals nor datum-dependent or time-dependent quadratic forms. An
argument requiring monotonicity only in a small-data region also lies
outside the theorem.

For a classical solution write
\begin{equation}
 E(t)=\norm{u(t)}_2^2,\quad Y(t)=\norm{\nabla u(t)}_2^2,
 \quad \omega=\curl u,\quad \Omega(t)=\norm{\omega(t)}_\infty.
 \label{eq:observables}
\end{equation}
The upper right Dini derivative is
$D^+\Omega(t)=\limsup_{h\downarrow0}(\Omega(t+h)-\Omega(t))/h$.

\begin{theorem}[Matched scalar values and unbounded initial growth]
\label{thm:growth}
For each $\nu>0$ there are constants $E_*,Y_*,\gamma>0$, a compact set
$K\subset\R^3$, and real odd fields
$u_{0,N}\in C_c^\infty(\R^3;\R^3)$, $N\ge1$, such that:
\begin{enumerate}
\item The data are solenoidal, supported in $K$, uniformly bounded in
$L^\infty$, and satisfy
\begin{equation}
 E(u_{0,N})=E_*,\qquad Y(u_{0,N})=Y_*,\qquad
 \norm{\curl u_{0,N}}_\infty=1.\label{eq:matched}
\end{equation}
\item For a radius $\varepsilon_N>0$,
$\curl u_{0,N}=e_3$ on $B_{\varepsilon_N}(0)$. In particular, the origin is a
global vorticity maximizer and all positive-order spatial derivatives of
vorticity vanish there.
\item The solution of \eqref{eq:NS} from every $u_{0,N}$ is global and
smooth, and
\begin{equation}
 \liminf_{t\downarrow0}\frac{\Omega_N(t)-1}{t}\ge\gamma N.
 \label{eq:unbounded-growth}
\end{equation}
\end{enumerate}
Moreover, the common product $E_*Y_*$ can be made arbitrarily small while
retaining these assertions.
\end{theorem}

Only the \emph{initial} supports are compact. No compact-support propagation
for the equation is asserted. The radii $\varepsilon_N$ shrink with $N$,
and \eqref{eq:unbounded-growth} contains no time interval uniform in $N$.
Thus the result is not a norm-inflation theorem and does not assert
persistence of large growth.

\subsection{Relation to existing work}
The large-amplitude obstruction to using general quadratic functionals as
global Lyapunov functions is established background. Goulart and
Chernyshenko \cite{GC} discuss why the highest-degree advective contribution
must cancel. Darrow, Carlson, and Goluskin \cite{DCG}, in developing quartic
Lyapunov functions, explicitly discuss the energy-only quadratic restriction
for typical incompressible systems. Theorem~\ref{thm:quadratic} gives a
precise whole-space statement with $H^m$ continuity and $L^3$ coercivity,
without imposing symmetries on the original form. It is not a claim to have
discovered the qualitative energy-only principle.

Quadratic invariants of Fourier-truncated inviscid dynamics were already
studied by Lee \cite{Lee} and in the helical equilibrium literature
\cite{Kraichnan}. Here a localized interaction supplies the necessary
scalar-symbol classification directly; no truncated-system classification
is imported as a theorem about the full equation. In particular, the
orthogonal averaging includes reflections, not just rotations.

Large instantaneous growth must also be distinguished from sustained
growth. Lu and Doering \cite{LD} study bounds and constrained optimization
of enstrophy production, with numerical evidence for the three-dimensional
extremizers. Kang, Yun, and Protas \cite{KYP} study finite-time enstrophy
amplification numerically. Theorem~\ref{thm:growth} instead fixes energy,
enstrophy, \emph{and} maximum vorticity exactly, and proves unbounded initial
growth of the last quantity analytically on globally regular whole-space
solutions. Nonlocal strain--vorticity interactions are familiar
\cite{BPB}; local directional coherence on a shrinking ball is not the
quantitative geometric control entering regularity criteria such as
\cite{CF}. The theorem makes this distinction through an explicit family.

Tao's averaged-equation construction \cite{TaoAveraged} provides a different
kind of obstruction: energy cancellation survives in a modified equation
that admits blowup. Neither result here modifies the equation. Conversely,
neither result establishes that every conceivable regularity proof must
use a particular combination of nonlocal, temporal, or full-datum
information. Their conclusions apply only to the quantified classes above.

\subsection{Organisation and scope}
Section~\ref{sec:prelim} records the local theory used and proves the
estimates needed below. Section~\ref{sec:quadratic} proves
Theorem~\ref{thm:quadratic}, including localization for measurable symbols.
Section~\ref{sec:growth} proves Theorem~\ref{thm:growth}.
Section~\ref{sec:consequences} states the exact comparison-law exclusions.
The proofs of both main results are contained in this manuscript; only
standard local Navier--Stokes theory and standard functional analysis are
used as background. This is an author-checked manuscript, not a claim of
independent expert certification or established priority for either exact
formulation.

\section{Preliminaries}\label{sec:prelim}
We use the unitary Fourier transform, with $\widehat{\partial_jv}(\xi)=
i\xi_j\widehat v(\xi)$. The Leray projection $\Pj$ has symbol
\begin{equation}
 P(\xi)=I-\frac{\xi\otimes\xi}{|\xi|^2},\qquad \xi\ne0.
 \label{eq:leray}
\end{equation}
It is bounded on each $H^s$ and self-adjoint on $L^2$. Put
$\N(w)=\Pj((w\cdot\nabla)w)$. Then the projected equation is
$u_t=\nu\Delta u-\N(u)$.

\subsection{Local theory used}
We use the classical maximal solution supplied by local $H^1$ theory:
for solenoidal Schwartz data there is $T_*\in(0,\infty]$, with a unique
$H^1$ mild branch, such that on each $[0,T]$, $T<T_*$, the velocity belongs
to $C^j([0,T];H^k)$ for all nonnegative integers $j,k$. If $T_*<\infty$,
its $H^1$ norm is unbounded as $t\uparrow T_*$. The normalised pressure
is determined by \eqref{eq:NS} and the Leray projection.

At unit viscosity these facts follow from Theorem~5.4 and Corollary~5.8
of Tao \cite{TaoLocal}, with the smooth-data regularity in the same local
theorem. The common maximal branch is obtained by uniqueness and gluing
local branches; the supremum of their existence times is $T_*$. Smooth
spatial bounds and the projected equation give continuity and time
differentiability in each finite Sobolev norm, increasing the spatial
index before differentiating. For general viscosity set
$u(x,t)=\nu v(x,\nu t)$, $p(x,t)=\nu^2q(x,\nu t)$; then $v$ solves the
unit-viscosity equation with datum $u_0/\nu$. This establishes the version
used here without an inviscid limit. In particular differentiation at
$t=0$ in $H^m$ is legitimate for every datum used below.

\subsection{Energy and enstrophy estimates}
Let $Z=\norm{\Delta u}_2^2$ and
$S=\int\omega\cdot((\omega\cdot\nabla)u)$. On each compact classical
interval, integration by parts gives
\begin{equation}
 E'=-2\nu Y,\qquad \frac12Y'+\nu Z=S.\label{eq:balances}
\end{equation}
Indeed the velocity transport and pressure vanish against $u$, and the
vorticity equation
\begin{equation}
 \partial_t\omega+(u\cdot\nabla)\omega
       =(\omega\cdot\nabla)u+\nu\Delta\omega\label{eq:vorticity}
\end{equation}
paired with $\omega$ gives the second identity. Solenoidality and
Plancherel give $\norm\omega_2^2=Y$ and
$\norm{\nabla\omega}_2^2=Z$. Smooth Sobolev bounds justify these integrations,
for instance by cutoffs followed by their removal.

Sobolev interpolation yields
\begin{equation}
 |S|\le \norm\omega_3^2\norm{\nabla u}_3
       \le C Y^{3/4}Z^{3/4}.
\end{equation}
Here $\norm\omega_6\le C\norm{\nabla\omega}_2$ and
$\norm{\nabla u}_6\le C\norm{\nabla^2u}_2=C Z^{1/2}$.
Young's inequality and \eqref{eq:balances} imply, with an absolute $C_0>0$,
\begin{equation}
 Y'+\nu Z\le C_0\nu^{-3}Y^3.\label{eq:enstrophy-inequality}
\end{equation}
Also Fourier Cauchy--Schwarz gives
\begin{equation}
 Y^2\le EZ.\label{eq:EYCS}
\end{equation}

\begin{lemma}[A small-product region]\label{lem:small}
If a nonzero solenoidal Schwartz datum satisfies
$C_0E(0)Y(0)/\nu^4<1/2$, its classical solution is global and
$Y(t)\le Y(0)$ for all $t\ge0$.
\end{lemma}
\begin{proof}
Write $E_0=E(0)>0$ and $Y_0=Y(0)>0$. By
\eqref{eq:balances}--\eqref{eq:EYCS},
\begin{equation}
 Y'\le -\frac{\nu}{E_0}Y^2+C_0\nu^{-3}Y^3
       =-\frac{\nu}{E_0}Y^2
          \left(1-\frac{C_0E_0Y}{\nu^4}\right).\label{eq:barrier}
\end{equation}
At every state with $Y=Y_0$ the right side is strictly negative. Since
$Y$ is continuously differentiable and starts at $Y_0$, it initially
decreases and cannot subsequently cross that level upwards: at a first
return from below its derivative would be nonnegative. Thus $Y\le Y_0$
throughout the maximal interval. Together with $E\le E_0$ this rules out
the finite-time $H^1$ blowup alternative. The local theory gives global
smoothness on every finite interval.
\end{proof}

We will also use the elementary inequalities
\begin{equation}
 |S|\le\Omega\norm\omega_2\norm{\nabla u}_2=\Omega Y,
 \qquad \norm u_3^4\le C E Y.\label{eq:comparison}
\end{equation}
The second follows by interpolation between $L^2$ and $L^6$, then Sobolev.

\section{Proof of the quadratic obstruction}\label{sec:quadratic}
Assume that $B$ satisfies Theorem~\ref{thm:quadratic}'s two requirements.
We derive a contradiction.

\subsection{Amplitude testing forces an exact cancellation}
For $w\in C_c^\infty(\R^3;\R^3)$ solenoidal and $a\in\R$, the solution
from $aw$ has initial derivative $\nu a\Delta w-a^2\N(w)$. Thus
\begin{equation}
 \left.\frac{\mathrm d}{\mathrm dt}B(u(t),u(t))\right|_{t=0}
 =2\nu a^2B(w,\Delta w)-2a^3B(w,\N(w))\le0.\label{eq:amplitude}
\end{equation}
The vectors in these pairings lie in $H^m_\sigma$. Taking arbitrarily large
positive and negative $a$ forces
\begin{equation}
                 B(w,\N(w))=0
 \quad(w\in C_c^\infty,\ \diver w=0).\label{eq:cancellation}
\end{equation}
We have varied the initial datum. We have not used velocity sign reversal
as a symmetry of the solution equation.

\subsection{Translation and orthogonal averaging}
For $\tau_yv(x)=v(x-y)$ and $R>0$ define
\begin{equation}
 B_R(v,w)=\frac{1}{|B_R(0)|}\int_{B_R(0)}
                         B(\tau_yv,\tau_yw)\dd y.
\end{equation}
Translations are strongly continuous isometries on $H^m$ and $L^3$.
The forms $B_R$ are symmetric, obey the same continuity bound, and satisfy
\begin{equation}
 B_R(v,v)\ge c\norm v_3^2,
 \qquad B_R(w,\N(w))=0\quad(w\in C_c^\infty{}_\sigma).
 \label{eq:average-properties}
\end{equation}
For the cancellation use $\N(\tau_yw)=\tau_y\N(w)$ and
\eqref{eq:cancellation}. Choose a countable dense subset of $H^m_\sigma$.
Boundedness and a diagonal subsequence along $R\to\infty$ give convergence
on all pairs from that subset. The uniform continuity bound extends the
limiting bilinear form to all of $H^m_\sigma$ and gives convergence on
every fixed pair there. Denote this form by $B_\infty$; it retains
\eqref{eq:average-properties}.

For fixed $z\in\R^3$,
\begin{align}
 &|B_R(\tau_zv,\tau_zw)-B_R(v,w)|\notag\\
 &\hspace{1cm}\le C\norm v_{H^m}\norm w_{H^m}
 \frac{|(B_R(0)+z)\mathbin\triangle B_R(0)|}{|B_R(0)|}
 \longrightarrow0.
\end{align}
Consequently $B_\infty$ is translation invariant. This compactness argument
concerns uniformly bounded bilinear forms, not solution trajectories.

For $O\in O(3)$ write $T_Ov(x)=O v(O^Tx)$ and use probability Haar measure
to set
\begin{equation}
 \overline B(v,w)=\int_{O(3)}B_\infty(T_Ov,T_Ow)\dd O.
 \label{eq:orthogonal-average}
\end{equation}
Orthogonal changes of variables preserve the relevant norms,
solenoidality, and covariance of $\N$, including when $\det O=-1$.
Thus $\overline B$ is symmetric, translation invariant and $O(3)$ invariant,
and retains coercivity, continuity, and cancellation. We henceforth work
with $\overline B$.

\subsection{The averaged multiplier}
\begin{lemma}\label{lem:multiplier}
There is a real measurable function $q:(0,\infty)\to\R$ such that
\begin{equation}
 |q(r)|\le C(1+r^2)^m\quad\text{for a.e. }r>0,
 \qquad
 \overline B(v,v)=\int_{\R^3}q(|\xi|)|\widehat v(\xi)|^2\dd\xi
 \quad(v\in H^m_\sigma).
 \label{eq:radial-multiplier}
\end{equation}
\end{lemma}
\begin{proof}
Extend the form to all real vector-valued $H^m$ by inserting $\Pj$ in
both slots. Let $J=(1-\Delta)^{m/2}:H^m\to L^2$. The form
$\overline B(\Pj J^{-1}f,\Pj J^{-1}g)$ is represented on real $L^2$ by a
bounded self-adjoint operator. Complexify this operator. It commutes with
translations, so its Fourier transform commutes with multiplication by
every character $e^{iy\cdot\xi}$. By the spectral theorem for the three
coordinate multiplication operators, it commutes with their spectral
projections, and therefore with every bounded scalar multiplication
operator.

For clarity, the last commutation property implies matrix multiplication
without any regularity assumption on the symbol. Exhaust frequency space
by nested finite-measure sets $K_n$. Apply the operator to
$1_{K_n}e_j$. Commutation with $1_{K_n}$ shows that the results for larger
sets restrict consistently to $K_n$, defining measurable columns of a
matrix $A(\xi)$. Commutation with indicators then identifies the operator
on all simple vector functions supported in $K_n$. For any constant vector
$a$ and measurable $F\subset K_n$, its norm bound gives
$\int_F|A(\xi)a|^2\dd\xi\le C^2|a|^2|F|$.
A countable dense set of $a$ implies $\norm{A(\xi)}\le C$ a.e.; density
identifies the operator on all $L^2$. Self-adjointness makes $A$ Hermitian.

Transferring back through $J$ gives a measurable Hermitian symbol
$M(\xi)$ for the $H^m$ form, with
\begin{equation}
 M=PMP,\qquad \norm{M(\xi)}\le C(1+|\xi|^2)^m.\label{eq:matrix-bound}
\end{equation}
There is no distribution supported at $\xi=0$ in this representation.
To justify pointwise symmetry despite possible exceptional null sets,
one may first represent $B_\infty$ by its symbol $M_\infty$ and then form
\begin{equation}
 M_*(\xi)=\int_{O(3)}O^T M_\infty(O\xi)O\dd O.\label{eq:symbol-average}
\end{equation}
Fubini and polar coordinates make this integral well-defined on every
direction for almost every radius. It represents
\eqref{eq:orthogonal-average} and gives an equivariant representative:
$M_*(R\xi)=R M_*(\xi)R^T$. Indeed substitute $O'=OR$ in
\eqref{eq:symbol-average}. The same construction retains
\eqref{eq:matrix-bound}.

At a fixed nonzero $\xi$, its stabilizer in $O(3)$ acts as the full $O(2)$
on $\xi^\perp$. Commutation with reflections about the two coordinate
axes of this plane eliminates off-diagonal entries of its Hermitian
restriction. Interchange of these axes makes the diagonal entries equal.
The longitudinal restriction is zero by $M=PMP$. Hence
$M_*(\xi)=q(|\xi|)P(\xi)$, with $q$ real; equivariance makes it radial.
This proves \eqref{eq:radial-multiplier}. Full orthogonal averaging is
essential: averaging only over $SO(3)$ would leave a possible helical
component on the transverse plane.
\end{proof}

\subsection{Localization with a measurable radial weight}
We need a localization statement because an infinite periodic field is not
an admissible datum on $\R^3$.

\begin{lemma}\label{lem:localization}
Let $q$ satisfy the bound in \eqref{eq:radial-multiplier}, and let
$\chi\in C_c^\infty(\R^3)$ be radial. There is a full-measure set
$G\subset(0,\infty)$ such that, for every nonzero $\eta$ with $|\eta|\in G$
and every $a\in\C^3$,
\begin{equation}
 \norm{q(D)(\chi(x/L)a e^{i\eta\cdot x})
      -q(|\eta|)\chi(x/L)a e^{i\eta\cdot x}}_2=o(L^{3/2})
 \quad(L\to\infty).
 \label{eq:localization}
\end{equation}
Here $q(D)$ denotes the scalar Fourier multiplier with symbol $q(|\xi|)$.
\end{lemma}
\begin{proof}
Choose $G$ to consist of the positive Lebesgue points of both $q$ and
$q^2$, with a compatible representative of $q$. For $r_0\in G$ the
one-dimensional mean of $|q(r)-q(r_0)|^2$ near $r_0$ tends to zero.
If $|\eta|=r_0$ and $0<\varepsilon<r_0/2$, the area of each spherical
slice of $B_\varepsilon(\eta)$ is bounded by $C\varepsilon^2$.
The coarea formula therefore gives
\begin{equation}
 \frac{1}{|B_\varepsilon|}
 \int_{B_\varepsilon(\eta)}|q(|\xi|)-q(r_0)|^2\dd\xi
 \le\frac{C}{\varepsilon}
       \int_{r_0-\varepsilon}^{r_0+\varepsilon}
                       |q(r)-q(r_0)|^2\dd r\longrightarrow0.
 \label{eq:radial-lebesgue}
\end{equation}
This works in every direction at radius $r_0$.

Plancherel and the change $\xi=\eta+\zeta/L$ show that the square of the
left side of \eqref{eq:localization}, divided by $L^3$, equals
\begin{equation}
 |a|^2\int_{\R^3}|q(|\eta+\zeta/L|)-q(r_0)|^2
                              |\widehat\chi(\zeta)|^2\dd\zeta.
 \label{eq:localization-integral}
\end{equation}
Near $\eta$, this tends to zero by \eqref{eq:radial-lebesgue} and the
approximate-identity property of a rapidly decaying radial kernel.
One can verify that property directly: dominate $|\widehat\chi(\zeta)|^2$
by $C_M(1+|\zeta|)^{-M}$, split into dyadic annuli, and bound each annular
integral by the mean over the containing ball. For balls of sufficiently
small radius those means are as small as desired; the remaining annular
weights have a summable tail. Local boundedness of $q$ controls them.
For $|\zeta|>\delta L$, with fixed $\delta>0$, the polynomial bound on $q$
gives a dominating constant times
$(1+|\zeta|)^{4m}|\widehat\chi(\zeta)|^2$, whose tail integral tends to
zero. This proves \eqref{eq:localization-integral} tends to zero without
assuming continuity of $q$.
\end{proof}

\subsection{A three-component interaction forces a constant weight}
Fix $k,l>0$, set $r=(k^2+l^2)^{1/2}$, and consider the finite Fourier fields
\begin{align}
 U&=(0,-\cos(kx_1),\ \sin(lx_2)+\cos(kx_1+lx_2)),\label{eq:triad-U}\\
 A&=(l^{-1}\cos(lx_2),\ k^{-1}\sin(kx_1+lx_2),\ k^{-1}\sin(kx_1)).
 \label{eq:triad-A}
\end{align}
Direct differentiation gives $\curl A=U$ and $\diver U=0$, while
\begin{equation}
 (U\cdot\nabla)U
  =(0,0,-l\cos(kx_1)[\cos(lx_2)-\sin(kx_1+lx_2)]).
 \label{eq:triad-nonlinearity}
\end{equation}
Temporarily interpreting $q(D)$ on these finitely many modes,
\begin{equation}
 \mean\{q(D)U\cdot (U\cdot\nabla)U\}
                     =\frac l4\{q(l)-q(r)\}.\label{eq:triad-mean}
\end{equation}
The mean is over the two periods in $x_1,x_2$; the fields are independent
of $x_3$. The two surviving products use
$\mean(\sin(lx_2)\cos(kx_1)\sin(kx_1+lx_2))=1/4$ and
$\mean(\cos(kx_1+lx_2)\cos(kx_1)\cos(lx_2))=1/4$.

Choose a nonnegative nonzero radial $\chi\in C_c^\infty$ and put
$\chi_L(x)=\chi(x/L)$,
\begin{equation}
 W_L=\curl(\chi_L A)=\chi_LU+R_L,
 \qquad R_L=\nabla\chi_L\times A.\label{eq:localized-triad}
\end{equation}
Now $W_L$ is a genuinely compactly supported smooth solenoidal field.
For every fixed integer $s\ge0$,
\begin{align}
 \norm{R_L}_{H^s}&=O(L^{1/2}),&
 \norm{R_L}_{W^{1,\infty}}&=O(L^{-1}),\label{eq:remainder-est}\\
 \norm{W_L\cdot\nabla W_L-\chi_L^2U\cdot\nabla U}_2
                         &=O(L^{1/2}).\label{eq:nonlinear-est}
\end{align}
To see these bounds, every cutoff derivative supplies $L^{-1}$,
all derivatives of $U,A$ are bounded, and the supports have volume
$O(L^3)$. In \eqref{eq:nonlinear-est} the remaining main error is
$\chi_L(U\cdot\nabla\chi_L)U$, with the same estimates.

Choose $k,l,r\in G$, where $G$ is furnished by
Lemma~\ref{lem:localization}. Apply that lemma to each of the finite
frequencies of $U$. The symbol bound and \eqref{eq:remainder-est} at
$s=2m$ also give $\norm{q(D)R_L}_2=O(L^{1/2})$. Consequently
\begin{equation}
 q(D)W_L=\chi_Lq(D)U+o_{L^2}(L^{3/2}).\label{eq:main-localized}
\end{equation}
The scalar multiplier preserves solenoidality. By polarization of
\eqref{eq:radial-multiplier}, moving the multiplier to the smooth first
factor and removing the Leray projection is legitimate. Thus
\begin{align}
 0&=\overline B(W_L,\N(W_L))\notag\\
  &=\ip{q(D)W_L}{W_L\cdot\nabla W_L}\notag\\
  &=\int\chi_L^3\,q(D)U\cdot(U\cdot\nabla U)+o(L^3)\notag\\
  &=\frac l4\{q(l)-q(r)\}L^3\int\chi^3+o(L^3).
 \label{eq:triad-conclusion}
\end{align}
For the last equality, each nonconstant Fourier mode has integral
$O(L^2)$ after one integration by parts against $\chi_L^3$; the constant
mode is \eqref{eq:triad-mean}. Hence $q(l)=q(r)$.

The conditions $k,l,(k^2+l^2)^{1/2}\in G$ hold for almost every pair
$(k,l)\in(0,\infty)^2$. The smooth change of variables
$(k,l)\mapsto(r,l)$ is a local diffeomorphism onto $0<l<r$, so
$q(l)=q(r)$ for almost every such pair. Symmetry and Fubini imply that
$q$ is constant a.e. on every finite interval separated from zero;
overlap makes the constant the same on $(0,\infty)$. Denote it by $a_0$.

The third component in \eqref{eq:triad-U} is important. This argument does
not classify quadratic invariants of two-component two-dimensional
velocity fields, which possess enstrophy as well as energy.

\subsection{Contradiction with critical coercivity}
The preceding argument gives $\overline B(v,v)=a_0\norm v_2^2$. But its
coercivity was retained during both averages. For a nonzero
$v\in C_c^\infty{}_\sigma$ set $v_\lambda(x)=\lambda v(\lambda x)$. Then
\begin{equation}
 c\norm v_3^2=c\norm{v_\lambda}_3^2
 \le a_0\norm{v_\lambda}_2^2
 =a_0\lambda^{-1}\norm v_2^2\longrightarrow0
 \quad(\lambda\to\infty),
\end{equation}
a contradiction. This proves Theorem~\ref{thm:quadratic}.
We classified only the averaged scalar symbol; no classification of the
original unsymmetrized bilinear form is needed or asserted.

\section{A globally smooth matched-scalar family}\label{sec:growth}
We now give the complete construction for Theorem~\ref{thm:growth}.
The nested fields produce strain at the origin while keeping their
vorticity supports disjoint. Two remote reservoirs then match energy and
enstrophy exactly. A small-product choice places every datum in
Lemma~\ref{lem:small}'s global region.

\subsection{Compact strain with annular vorticity}
Let
\begin{equation}
 H=\diag(-1/2,-1/2,1),\qquad
 A(x)=(-x_2x_3/2,x_1x_3/2,0).
\end{equation}
One has $\curl A=Hx$. Choose a smooth radial cutoff $\chi$ equal to one
on $B_1$ and zero outside $B_2$, and set $W=\curl(\chi A)$. It is odd,
compactly supported, and solenoidal. On $B_1$ it equals $Hx$, which is
irrotational. Therefore $\curl W$ is supported in
$\{1\le|x|\le2\}$. Put
\begin{equation}
 C=\norm{\curl W}_\infty>0,\qquad \gamma=\frac{1}{4C}.
 \label{eq:gamma}
\end{equation}
If $C$ were zero, the compact field would have zero curl and divergence;
Plancherel would force it to vanish, contrary to $W=Hx$ on $B_1$.
For $r>0$ define
\begin{equation}
 W_r(x)=rW(x/r).\label{eq:strain-scaling}
\end{equation}
It has gradient $H$ on $B_r$, and its vorticity has supremum $C$ and is
supported in $\{r\le|x|\le2r\}$. Also
$E(W_r)=r^5E(W)$, $Y(W_r)=r^3Y(W)$, and
$\norm{W_r}_\infty=r\norm W_\infty$.

\subsection{A rotation core with a certified global maximum}
Choose $g\in C_c^\infty((0,4))$ with $0\le g\le1$ and
$\int g=1$. Define
\begin{equation}
 f(r)=\begin{cases}
  1,&0\le r\le1,\\
  1-\displaystyle\int_0^{\log r}g(s)\dd s,&r>1,
 \end{cases}
 \qquad R=e^4.
\end{equation}
Then $f=0$ for $r\ge R$, $0\le f\le1$, and $-1\le rf'(r)\le0$;
the joins are smooth. Put
\begin{equation}
 V(x)=\tfrac12 f(|x|)(-x_2,x_1,0).\label{eq:rotation}
\end{equation}
It is odd, smooth, compactly supported and solenoidal, with curl $e_3$ on
$B_1$. For $n=x/|x|$ and $c=n_3$ direct differentiation gives
\begin{align}
 \curl V&=(f+rf'/2)e_3-(rf'/2)c n,\notag\\
 |\curl V|^2&=f^2c^2+(f+rf'/2)^2(1-c^2)\le1.\label{eq:rotation-bound}
\end{align}
The inequality uses $f\in[0,1]$ and $f+rf'/2\in[-1/2,1]$.
Thus $\norm{\curl V}_\infty=1$ exactly, including the cutoff region.
For $\varepsilon>0$ set
$V_\varepsilon(x)=\varepsilon V(x/\varepsilon)$. Its vorticity has the
same supremum and equals $e_3$ on $B_\varepsilon$.

\subsection{Nested fields and uniform bounds}
Choose $r_0\in(0,1)$ later, and put
\begin{equation}
 r_k=r_0 4^{-(k-1)},\quad
 \varepsilon_N=\frac{r_N}{4R},\quad
 C_N=\gamma\sum_{k=1}^N W_{r_k}+V_{\varepsilon_N}.
 \label{eq:central}
\end{equation}
The vorticity annuli of the $W_{r_k}$ are mutually disjoint.
The rotation core is supported in $B_{r_N/4}$, disjoint from every annulus.
By \eqref{eq:gamma} and \eqref{eq:rotation-bound},
\begin{equation}
 \norm{\curl C_N}_\infty=1.
\end{equation}
On $B_{\varepsilon_N}$,
\begin{equation}
 C_N(x)=(\gamma N H+K_0)x,\qquad
 K_0=\begin{pmatrix}0&-1/2&0\\1/2&0&0\\0&0&0\end{pmatrix}.
 \label{eq:affine-core}
\end{equation}
In particular $C_N(0)=0$ and $\curl C_N=e_3$ near zero.

The velocity supports of the nested strain fields overlap, so energy
must not be added as if they were disjoint. Instead the triangle inequality
and scaling give
\begin{equation}
 \norm{C_N}_2\le
 \frac{\gamma\norm W_2 r_0^{5/2}}{1-4^{-5/2}}
       +\norm V_2\left(\frac{r_0}{4R}\right)^{5/2}.
 \label{eq:central-energy}
\end{equation}
For enstrophy the vorticity supports \emph{are} disjoint, and the div--curl
identity gives the exact formula
\begin{equation}
 Y(C_N)=\gamma^2Y(W)\sum_{k=1}^N r_k^3
                         +\varepsilon_N^3Y(V).
 \label{eq:central-enstrophy}
\end{equation}
Hence there are constants $C_E,C_Y$ independent of $N,r_0$ such that
\begin{equation}
 E(C_N)\le C_Er_0^5,\qquad Y(C_N)\le C_Yr_0^3.
 \label{eq:central-small}
\end{equation}
The analogous convergent sum of $r_k$ bounds $\norm{C_N}_\infty$
uniformly in $N$. Every $C_N$ is supported in $B_{2r_0}$.

\subsection{Exact matching with two remote reservoirs}
Take fixed length scales $L_1=1,L_2=2$. Choose fixed centres
$\pm a_1,\pm a_2$ so that the four closed balls of radii $RL_j$ are
pairwise disjoint and avoid $\overline B_2$. For a parameter
$\delta\in(0,1/8)$ put
\begin{equation}
 R_j(x)=\delta L_j
 \left[V\left(\frac{x-a_j}{L_j}\right)
       +V\left(\frac{x+a_j}{L_j}\right)\right],\quad j=1,2.
 \label{eq:reservoirs}
\end{equation}
They are odd and solenoidal, have disjoint supports, and satisfy
\begin{equation}
 E_j:=E(R_j)=2\delta^2L_j^5 E(V),\quad
 Y_j:=Y(R_j)=2\delta^2L_j^3 Y(V),\quad
 \norm{\curl R_j}_\infty=\delta.
 \label{eq:reservoir-scalars}
\end{equation}
Let
\begin{equation}
 M=\begin{pmatrix}E_1&E_2\\Y_1&Y_2\end{pmatrix},\qquad
 E_*=E_1+E_2,\qquad Y_*=Y_1+Y_2.
\end{equation}
The matrix is invertible, since
$Y_j/E_j=L_j^{-2}Y(V)/E(V)$ differs for $j=1,2$.
Moreover
\begin{equation}
 E_*Y_*=C_R\delta^4\label{eq:reservoir-small}
\end{equation}
for a positive constant $C_R$ determined by $V,L_1,L_2$.

First choose $\delta>0$ sufficiently small that
$C_0E_*Y_*/\nu^4<1/2$. The product in
\eqref{eq:reservoir-small} may also be made smaller than any prescribed
positive number. Next choose $r_0>0$ sufficiently small and define
\begin{equation}
 \binom{a_N^2}{b_N^2}
   =\binom11-M^{-1}\binom{E(C_N)}{Y(C_N)}.
 \label{eq:matching-coefficients}
\end{equation}
Indeed $M=\delta^2 M_0$ with $M_0$ fixed and invertible; by
\eqref{eq:central-small} one can choose $r_0$, after $\delta$, such that
\begin{equation}
                 \frac34\le a_N^2,b_N^2\le\frac54
                 \qquad\text{for every }N.\label{eq:coefficient-bounds}
\end{equation}
Take the positive square roots and set
\begin{equation}
                   u_{0,N}=C_N+a_NR_1+b_NR_2.\label{eq:final-data}
\end{equation}
The three terms in \eqref{eq:final-data} now have disjoint
\emph{velocity} supports. Equations \eqref{eq:matching-coefficients} and
\eqref{eq:reservoir-scalars} give $E(u_{0,N})=E_*$ and
$Y(u_{0,N})=Y_*$ exactly. The reservoirs have vorticity at most
$\sqrt{5/4}\,\delta<1$, so the exact maximum vorticity remains one.
They vanish near the origin and change none of \eqref{eq:affine-core}.
Oddness, common compact initial support, and the uniform velocity bound
follow from their construction and \eqref{eq:coefficient-bounds}.

\subsection{Global smoothness and actual initial growth}
For the fixed viscosity $\nu$, Lemma~\ref{lem:small} applies to every datum
\eqref{eq:final-data}, with the \emph{same} constants $E_*,Y_*$. Thus each
solution is globally smooth and obeys $E_N(t)\le E_*$, $Y_N(t)\le Y_*$.
This uses only smallness of their common product, not arbitrary-data global
regularity.

Evaluate the actual vorticity equation \eqref{eq:vorticity} at $(0,0)$.
The advection term vanishes because the vorticity is constant near zero;
its Laplacian vanishes for the same reason. Equations
\eqref{eq:affine-core} and \eqref{eq:final-data} yield
\begin{equation}
 \partial_t\omega_N(0,0)
       =(e_3\cdot\nabla)u_{0,N}(0)=\gamma N e_3.
 \label{eq:actual-generator}
\end{equation}
Local smoothness up to $t=0$ gives
\begin{equation}
 |\omega_N(0,t)|=1+\gamma Nt+o_N(t).
\end{equation}
Since $\Omega_N(t)\ge|\omega_N(0,t)|$ and $\Omega_N(0)=1$, this proves
\eqref{eq:unbounded-growth} and completes Theorem~\ref{thm:growth}.
No differentiability of the maximum or persistence of its spatial location
has been assumed. The error $o_N(t)$ is not asserted uniform in $N$.

\section{Consequences and limitations}\label{sec:consequences}
\begin{corollary}[No finite scalar instantaneous bound]\label{cor:pointwise}
There is no function $F$, finite at every positive scalar tuple, such that
\begin{equation}
 D^+\Omega(t)\le F(\nu,E(t),Y(t),\Omega(t))\label{eq:scalar-law}
\end{equation}
holds at every classical state, including $t=0$, of every global smooth
solution from solenoidal Schwartz data. Adding the initial scalar values
$E(0),Y(0),\Omega(0)$ as arguments does not change this conclusion.
\end{corollary}
\begin{proof}
At $t=0$, Theorem~\ref{thm:growth} gives the same input tuple
$(\nu,E_*,Y_*,1)$ for every $N$, but a lower bound $\gamma N$ for its upper
right derivative. The initial scalar values are identical as well.
\end{proof}

\begin{corollary}[An almost-everywhere version]\label{cor:ae}
No locally bounded function $F$ of the same scalar arguments can make
\eqref{eq:scalar-law} hold merely for almost every time on an initial
classical interval of every solution in Theorem~\ref{thm:growth}.
\end{corollary}
\begin{proof}
The scalar tuple is continuous in time on each branch. Local boundedness
provides a neighborhood of the common initial tuple and a finite upper
bound $B$ for $F$ on it. For each $N$, the tuple remains in that neighborhood
on a positive interval, possibly depending on $N$. The smooth-data local
theory bounds $\partial_t\omega_N$ in $L^\infty$ on compact time intervals;
therefore $\Omega_N$ is locally Lipschitz. At its a.e. differentiability
points the assumed inequality gives $\Omega_N'\le B$, so integration gives
$\Omega_N(t)\le1+Bt$ on that interval. This contradicts
\eqref{eq:unbounded-growth} when $\gamma N>B$.
\end{proof}

At the exhibited point every spatial derivative of $\omega_{0,N}$ of
positive order is zero, and its value is $e_3$. The stretching in that
direction nevertheless equals $\gamma N$. Thus no finite bound for the
\emph{pointwise initial vorticity growth at this maximizing point} can
be recovered from the same scalar inputs and any finite, or the entire,
spatial vorticity jet there. This assertion is not about the full velocity
jet: its strain is precisely the varying quantity. A rule using all
maximizers or remote geometry is a different premise class.

\subsection{Why the excluded comparisons would have been useful}
For perspective, suppose an input-controlled positive continuous
comparison function $g$ gave $D^+\Omega\le g(\Omega)$ and
$\int_1^\infty g(s)^{-1}\dd s=\infty$, so that its scalar comparison
solution remains finite on finite intervals. Then
\eqref{eq:balances} and \eqref{eq:comparison} imply
\begin{equation}
 Y(t)\le Y(0)\exp\left(2\int_0^t\Omega(s)\dd s\right).
\end{equation}
An input-derived finite bound for this integral would contradict the local
blowup alternative at any finite maximal time. The corollaries exclude
only the indicated scalar-input instantaneous producers of this bound;
they do not exclude an estimate using additional controlled information
about the full initial datum or the evolving spatial field.

The first theorem has a similarly limited meaning. A fixed monotone form
coercive in $L^3$ would immediately bound that critical norm on the branch.
Theorem~\ref{thm:quadratic} prevents this particular universal quadratic
construction; it does not prevent nonlinear, time-dependent, or
datum-adapted constructions, nonmonotone estimates, or conditional
small-data results such as Lemma~\ref{lem:small}.

\subsection{No conclusion about arbitrary-data breakdown}
The coherent ball in Theorem~\ref{thm:growth} shrinks as $N$ grows. High
initial Sobolev norms need not be uniformly bounded. A large initial
vorticity derivative is compatible with global regularity, as the same
construction proves. These facts prevent the initial derivative from being
silently converted into a uniform-duration growth statement. The two main
theorems neither resolve arbitrary-data global regularity nor establish
that its remaining obstruction is strictly weaker. They delimit two
specified classes of proof mechanisms and leave other classes untouched.

\begin{thebibliography}{99}
\bibitem{BPB}
D.~Buaria, A.~Pumir, and E.~Bodenschatz,
\emph{Self-attenuation of extreme events in Navier--Stokes turbulence},
Nature Communications \textbf{11} (2020), 5852.
\href{https://doi.org/10.1038/s41467-020-19530-1}{doi:10.1038/s41467-020-19530-1}.

\bibitem{CF}
P.~Constantin and C.~Fefferman,
\emph{Direction of vorticity and the problem of global regularity for the
Navier--Stokes equations}, Indiana Univ. Math. J. \textbf{42} (1993), 775--789.
\href{https://doi.org/10.1512/iumj.1993.42.42034}{doi:10.1512/iumj.1993.42.42034}.

\bibitem{DCG}
D.~Darrow, E.~Carlson, and D.~Goluskin,
\emph{Quartic Lyapunov functions for global fluid stability},
arXiv:2606.18232v1 (16 June 2026).
\url{https://arxiv.org/abs/2606.18232}.

\bibitem{GC}
P.~J.~Goulart and S.~I.~Chernyshenko,
\emph{Global stability analysis of fluid flows using sum-of-squares},
Physica D \textbf{241} (2012), 692--704.
\href{https://doi.org/10.1016/j.physd.2011.12.008}{doi:10.1016/j.physd.2011.12.008}.

\bibitem{KYP}
D.~Kang, D.~Yun, and B.~Protas,
\emph{Maximum amplification of enstrophy in three-dimensional
Navier--Stokes flows}, J. Fluid Mech. \textbf{893} (2020), A22.
\url{https://arxiv.org/abs/1909.00041}.

\bibitem{Kraichnan}
R.~H.~Kraichnan,
\emph{Helical turbulence and absolute equilibrium},
J. Fluid Mech. \textbf{59} (1973), 745--752.
\href{https://doi.org/10.1017/S0022112073001837}{doi:10.1017/S0022112073001837}.

\bibitem{Lee}
J.~Lee,
\emph{Isolating constants of motion for the homogeneous turbulence of two
and three dimensions}, J. Math. Phys. \textbf{16} (1975), 1367--1373.
\href{https://doi.org/10.1063/1.522705}{doi:10.1063/1.522705}.

\bibitem{LD}
L.~Lu and C.~R.~Doering,
\emph{Limits on enstrophy growth for solutions of the three-dimensional
Navier--Stokes equations}, Indiana Univ. Math. J. \textbf{57} (2008), 2693--2728.
\href{https://doi.org/10.1512/iumj.2008.57.3716}{doi:10.1512/iumj.2008.57.3716}.

\bibitem{TaoLocal}
T.~Tao,
\emph{Localisation and compactness properties of the Navier--Stokes global
regularity problem}, Analysis \& PDE \textbf{6} (2013), 25--107.
\href{https://doi.org/10.2140/apde.2013.6.25}{doi:10.2140/apde.2013.6.25}.

\bibitem{TaoAveraged}
T.~Tao,
\emph{Finite time blowup for an averaged three-dimensional Navier--Stokes
equation}, J. Amer. Math. Soc. \textbf{29} (2016), 601--674.
\url{https://arxiv.org/abs/1402.0290}.
\end{thebibliography}
\end{document}
